\chapter{The Raft Consensus Engine Implementation}
\label{ch:raft_implementation}

The heart of the fault-tolerant key-value store is the Raft consensus engine. This chapter details the implementation of the Raft protocol in Go, focusing on the concurrency primitives, state management, and the safety mechanisms that ensure a consistent replicated log.

Implementing Raft in Go requires careful coordination of goroutines, condition variables, and channels to manage cluster events without deadlocks. This chapter focuses on persistent term tracking, leader election timeouts, and background replicators that synchronize log entries across cluster nodes.

\section{The Logical Clock: Utilizing Terms}
In the absence of synchronized physical clocks, the implementation relies on \textbf{terms} to provide a sense of logical time. Terms allow the system to detect and discard obsolete data, ensuring that only information from the most recent "epoch" is considered valid.

\subsection{Persistent Term Management}
The \texttt{Raft} structure maintains the logical clock via the \texttt{currentTerm} and \texttt{votedFor} fields. To survive crashes, these fields are persisted to stable storage whenever they are modified.

\begin{itemize}
    \item \textbf{currentTerm}: An integer that increments monotonically. It is used to identify the most recent leader and to ignore requests from "stale" nodes.
    \item \textbf{votedFor}: Tracks the candidate ID that received this node's vote in the current term, preventing double-voting.
\end{itemize}

\subsection{Detecting Obsolete State}
The implementation uses a centralized method, \texttt{handleHigherTerm}, to process incoming RPC terms.
If any node (including a leader) discovers that its \texttt{currentTerm} is less than the term in an incoming request or reply, it must immediately update its term and revert to the \texttt{Follower} state.

This logic is critical for safety.
By resetting \texttt{votedFor} to -1 upon a term increase, the implementation ensures the node is eligible to participate in the new term's election. This transition is also recorded in the \texttt{raft\_current\_term} metric to provide real-time visibility into the cluster's logical time progression.

\section{Leader Election Mechanics}
Leader election is the process by which the cluster maintains progress. The implementation utilizes Go's concurrency primitives—specifically goroutines, channels, and timers—to manage state transitions and timeouts.

\subsection{Ticker Routines and Randomized Timeouts}
Each node runs a background \texttt{ticker} goroutine that waits on two primary timers: the \texttt{electionTimer} and the \texttt{heartBeatTimer}. 

To prevent "split vote" scenarios, the \texttt{electionTimer} uses a randomized duration. The \texttt{randomizedElectionTimeout} method calculates a value between the minimum bound (\texttt{electionTimeoutMin}) and a random offset:

\begin{verbatim}
func (rf *Raft) randomizedElectionTimeout() time.Duration {
    ms := rf.electionTimeoutMin.Milliseconds() + 
          rand.Int63()%rf.electionTimeoutRand.Milliseconds()
    return time.Duration(ms) * time.Millisecond
}
\end{verbatim}

As depicted in the ticker logic, if the \texttt{electionTimer} fires while the node is a Follower or Candidate, it triggers a transition to the Candidate state.

\subsection{State Transitions: The becomeCandidate Routine}
When a node's election timeout expires, it transitions through the \texttt{becomeCandidate} method.
This routine performs the following atomic operations:
\begin{enumerate}
    \item Increments the \texttt{currentTerm}.
    \item Updates the state to \texttt{StateCandidate} and modifies the \texttt{raft\_state} metric.
    \item Votes for itself by setting \texttt{votedFor = rf.me}.
    \item Persists the new state and resets the election timer to start a new period.
    \item Spawns asynchronous \texttt{requestVoteFromPeer} goroutines to solicit votes from the rest of the cluster.
\end{enumerate}

This "Request-and-Vote" cycle allows a candidate to become the leader immediately upon receiving a majority of votes, at which point it calls \texttt{becomeLeader} to begin sending heartbeats and managing the log. A candidate becomes the leader immediately upon receiving a majority of votes, at which point it calls \texttt{becomeLeader} to begin sending heartbeats and managing the log. If a candidate receives an \texttt{AppendEntries} RPC from a valid leader during this process, it recognizes the leader's authority and reverts to the Follower state.

\begin{figure}[htbp]
    \centering
    \begin{tikzpicture}[
        node distance=3.6cm,
        state/.style={
            draw, circle, thick, fill=white,
            text width=2.2cm, minimum height=1.6cm, 
            align=center, font=\small
        },
        arrow/.style={-{Stealth}, thick, shorten >=2pt, shorten <=2pt}
    ]
        % Nodes
        \node[state, fill=orange!10] (Follower) at (0, 0) {Follower};
        \node[state, fill=blue!10] (Candidate) at (5, 0) {Candidate};
        \node[state, fill=green!10] (Leader) at (2.5, -3.2) {Leader};

        % Connectors
        \draw[arrow] (Follower) -- node[above, font=\scriptsize, align=center] {Election Timeout\\Fires} (Candidate);
        
        \draw[arrow] (Candidate) -- node[above, font=\scriptsize, align=center] {Wins Election\\(Majority)} (Leader);
        
        \draw[arrow] (Leader) .. controls +(left:2.0) and +(down:1.5) .. node[left, font=\scriptsize, align=right] {Discovers New\\Leader or Term} (Follower);
        
        \draw[arrow] (Candidate) .. controls +(up:2.0) and +(up:2.0) .. node[above, font=\scriptsize, align=center] {Higher Term\\Discovered} (Follower);
    \end{tikzpicture}
    \caption{State transitions and election mechanics within the Raft consensus node implementation.}
    \label{fig:raft_states}
\end{figure}

As illustrated in Figure \ref{fig:raft_states}, the state machine alternates between Follower, Candidate, and Leader based on election timeouts and term changes, managing progress and node liveness.

\section{Log Replication Mechanism}
\label{sec:log_replication_impl}
Once a leader is established, its primary responsibility is to replicate its log to a majority of followers. The implementation in \texttt{raft.go} utilizes a decoupled, concurrent architecture where each peer is managed by a dedicated goroutine.

The leader maintains two critical state indices for each follower: \texttt{nextIndex} and \texttt{matchIndex}. When new entries are submitted via the \texttt{Start} method, the replicator wakes up and pushes the entries to the peers using an \texttt{AppendEntries} RPC. If a follower's log is inconsistent, the leader can step backward to resolve the conflict, bypassing multiple index entries to converge safely on the correct state.

\subsection{Dedicated Replicator Routines}
For every peer in the cluster, the leader spawns a \texttt{replicator} goroutine during its transition to the \texttt{Leader} state. This routine operates as a long-running background process that synchronizes the leader's log with a specific follower.

The \texttt{replicator} routine follows a wait-and-signal pattern using condition variables (\texttt{sync.Cond}):
\begin{itemize}
    \item \textbf{Idle State}: The routine waits on \texttt{rf.replicatorCond[peer]} if no replication is currently required.
    \item \textbf{Triggering Replication}: Replication is triggered either by a heartbeat timer (to maintain authority) or by the submission of new client commands via \texttt{rf.Start()}.
    \item \textbf{Needs Replication Check}: The routine only proceeds if the node is still the leader and the peer's \texttt{matchIndex} is lagging behind the leader's last log index.
\end{itemize}

\subsection{Managing Next and Match Indices}
The leader maintains two critical state variables for each follower to manage the replication window:
\begin{itemize}
    \item \textbf{\texttt{nextIndex[]}}: The index of the next log entry the leader will send to that follower. It is initialized to the leader's last log index + 1.
    \item \textbf{\texttt{matchIndex[]}}: The index of the highest log entry known to be replicated on that follower. It is initialized to 0 and increases monotonically.
\end{itemize}

When an \texttt{AppendEntries} RPC fails due to log inconsistency, the leader resolves the conflict by adjusting the \texttt{nextIndex} backward.
The implementation optimizes this process using \texttt{conflictIndex} and \texttt{conflictTerm} metadata returned by the follower, allowing the leader to bypass multiple entries rather than decrementing by one index at a time.

The relationship between these indices ensures that the leader can independently track the progress of every node in the cluster. Successful responses from a follower trigger the \texttt{advancePeerIndices} method, which moves the \texttt{matchIndex} forward and signals the potential for a new commitment.

\section{Ensuring Safety Invariants}
The safety of the replicated state machine depends on strict rules governing when an entry is considered "committed" and safe to apply to the application state.

To ensure that committed data is never overwritten by a future leader, the implementation applies two conditions: the leader must have replicated the entry on a majority of nodes, and the entry must have been created during the leader's current term. This prevents "old" entries from being improperly overwritten, ensuring the cluster makes safe, consistent progress.

\subsection{Quorum Commitment and the Commit Index}
The leader advances its \texttt{commitIndex} by examining the \texttt{matchIndex} values of all peers.
The \texttt{updateCommitIndex} routine implements this logic by searching for the highest index $N$ that satisfies two conditions:
\begin{enumerate}
    \item \textbf{Majority Requirement}: $N$ must be replicated on a majority of the nodes in the cluster.
    \item \textbf{Current Term Requirement}: The entry at index $N$ must have been created in the leader's current term.
\end{enumerate}

This commitment logic ensures that once an entry is finalized, it cannot be overwritten by a subsequent leader. The restriction to only commit entries from the current term by counting replicas is a vital safety invariant that prevents the "overwriting committed entries" edge case discussed in the theoretical background.

\subsection{Log Application Rules}
Once an entry is committed, it must be applied to the state machine in order. This is handled by a dedicated \texttt{applier} goroutine that monitors the gap between \texttt{lastApplied} and \texttt{commitIndex}.

The \texttt{applier} ensures:
\begin{itemize}
    \item \textbf{Strict Ordering}: Entries are applied sequentially from \texttt{lastApplied + 1} up to \texttt{commitIndex}.
    \item \textbf{Asynchronous Notification}: Applied results are sent to the RSM via the \texttt{applyCh}, carrying the \texttt{Command}, \texttt{Index}, and \texttt{Term}.
    \item \textbf{Snapshot Integration}: If the node is behind and receives a snapshot, the applier prioritizes installing the snapshot state before resuming log application.
\end{itemize}

These rules, combined with the metrics tracking (\texttt{raft\_commit\_index} and \texttt{raft\_last\_applied}), provide a robust mechanism for maintaining state consistency across the distributed system.